\begin{table*}[t]
\centering
\caption{Supporting-title completeness and reader accuracy on the released IRCoT HotpotQA \texttt{test\_subsampled} split.}
\label{tab:evidence_completeness_buckets_supp}
\small
\begin{threeparttable}
\begin{tabularx}{\textwidth}{>{\raggedright\arraybackslash}X l r r r r}
\toprule
Scope & Supporting-title bucket & $N$ & EM & F1 & Mean title recall \\
\midrule
All retrieval-only method-example observations & full title support & 2181 & \best{0.7240} & \best{0.8584} & 1.0000 \\
All retrieval-only method-example observations & partial title support & 780 & 0.3846 & 0.4744 & 0.5000 \\
All retrieval-only method-example observations & no title support & 39 & 0.0769 & 0.0872 & 0.0000 \\
\midrule
HyDE-style RAG only & full title support & 434 & \best{0.7396} & \best{0.8720} & 1.0000 \\
HyDE-style RAG only & partial title support & 60 & 0.3500 & 0.4472 & 0.5000 \\
HyDE-style RAG only & no title support & 6 & 0.0000 & 0.0000 & 0.0000 \\
\bottomrule
\end{tabularx}
\begin{tablenotes}
\footnotesize
\item Full title support means all annotated supporting titles are retrieved in the top-10 evidence set; partial title support means at least one but not all supporting titles are retrieved. The pooled rows are descriptive method-example observations across the six retrieval-only methods and therefore are not treated as independent examples.
\end{tablenotes}
\end{threeparttable}
\end{table*}
